\clearpage
\thispagestyle{empty}
\pdfbookmark[0]{Ein neues Instrument zur Messung von Intensitätsänderungen. Ausgabe und Quellen}{instrument-1837-sources}
\begin{center}\Large\textbf{Ein neues Instrument zur Messung von Intensitätsänderungen}\end{center}

Diese Abhandlung von Carl Friedrich Gauss erschien in den \textit{Resultaten aus den Beobachtungen des magnetischen Vereins} für 1837, herausgegeben von Gauss und Wilhelm Weber und 1838 in Göttingen veröffentlicht. Die KI-Übertragung des vollständigen Textes und aller fünf historischen Anmerkungen in modernes Deutsch mit KI-erstelltem Satz folgt der gesammelten deutschen Ausgabe in den \textit{Werken}, Band V (Göttingen, 1867), mit dem unnummerierten Anfang und S.~358--373. Die Messwerte stehen in den historischen Einheiten Meile, Fuß, Zoll und Pfund. Die Schwingungsdauer bezeichnet hier die Zeit zwischen aufeinanderfolgenden äußersten Stellungen.

\medskip\noindent
William Francis übersetzte die Abhandlung für Richard Taylors \textit{Scientific Memoirs}, Bd.~II (1841), Teil VI, Artikel VI, S.~252--266. Major Sabine und Professor Lloyd revidierten die Übersetzung. Taylors gesonderter redaktioneller Anhang reicht bis S.~267. André Koch Torres Assis gibt die Übersetzung von Francis in \textit{Wilhelm Weber's Main Works on Electrodynamics Translated into English}, Bd.~I (2021), Kap.~11, S.~123--134, mit jeweils gekennzeichneten Anmerkungen von Gauss, Taylor und Assis wieder. Zur deutschen Terminologie der Instrumente siehe Julia Bloemer und Benjamin Mirwald (Bearb.), Deutsches Museum, \textit{Messingstab mit Magnet}, Nr.~1125, Fassung vom 24. August 2016.

\medskip\noindent
\href{https://gdz.sub.uni-goettingen.de/id/PPN236006339}{Gesammelte deutsche Ausgabe (1867)}\par
\noindent\href{https://archive.org/details/s1id11857360}{Übersetzung von Francis in den \textit{Scientific Memoirs} (1841)}\par
\noindent\href{https://www.ifi.unicamp.br/~assis/Weber-in-English-Vol-1.pdf\#page=126}{Ausgabe der Übersetzung von Francis durch Assis (2021)}\par
\noindent\href{https://digital.deutsches-museum.de/projekte/gruendungssammlung/detail/1125/}{Deutsches Museum, Instrumentenkatalog}
\clearpage
